\begin{definition}[Ultrafilter Lemma]
\label{def:ultrafilter-lemma}
The \emph{Ultrafilter Lemma} states the following.  Let \(I\) be a set and let
\(\mathcal F\subseteq\Pow(I)\) satisfy
\[
 I\in\mathcal F,\quad\varnothing\notin\mathcal F,\quad
 X,Y\in\mathcal F\Rightarrow X\cap Y\in\mathcal F,
\]
and
\[
 X\in\mathcal F\text{ and }X\subseteq Y\subseteq I
 \Rightarrow Y\in\mathcal F.
\]
Then there is \(\mathcal U\subseteq\Pow(I)\) containing \(\mathcal F\) that
satisfies the same conditions and, for every \(X\subseteq I\), contains
exactly one of \(X\) and \(I\setminus X\).  This is a metatheoretic choice
principle about families of subsets, not a formula of a first-order theory
under study.  It is strictly weaker over \(\ZF\) than the Axiom of Choice,
relative to the consistency of \(\ZF\) \cite{HalpernLevy1971}, and
therefore is not a second name for that axiom; its statused theorem form is
Theorem~\ref{thm:p01-w02-p01-009-ultrafilter-lemma}.  The later filter chapter gives
the short technical terms for the two displayed families before applying this
principle.
\end{definition}

Some model-theoretic constructions use the Ultrafilter Lemma, and some cardinal
arguments use the Axiom of Choice.  Later theorems state the required principle
rather than silently treating either principle as part of pure first-order
logic.

